\section{Теория Артина--Шрайера и векторы Витта}

Теория Куммера (\rf{5.27}--\rf{5.32}) описывает циклические расширения степени $n$ и абелевы расширения экспоненты, делящей $n$,
поля, содержащего первообразный корень степени $n$ из единицы. Если $\operatorname{char}F=p>0$ и $p\mid n$, такого корня
нет ни в каком поле характеристики $p$ (см.\ начало раздела <<Круговые расширения>>, перед \rf{5.7b}). В частности,
$X^p-1=(X-1)^p$, т.\,е.\ $\mu_p=\{1\}$. Присоединение корня $p$-й степени тоже бесполезно: если $a\notin F^p$,
то многочлен $X^p-a$ неприводим (\rf{2.3}), но в поле разложения $X^p-a=(X-\beta)^p$, так что $F[\beta]/F$~--- чисто
несепарабельное расширение степени $p$, и у него нет нетождественных $F$-автоморфизмов. Значит, циклические расширения
степени $p$ в характеристике $p$ устроены иначе (ср.\ \rf{5.29}(b)).

Выход~--- заменить мультипликативный гомоморфизм $x\mapsto x^n$ с ядром $\mu_n$ аддитивным гомоморфизмом
$\wp(x)=x^p-x$ с ядром $\F_p$, а мультипликативную теорему Гильберта~90~--- аддитивной. Так получается теория
Артина--Шрайера (расширения степени $p$), а для степеней $p^n$~--- теория Артина--Шрайера--Витта, в которой роль поля
играет кольцо векторов Витта длины $n$. Всюду в разделе $\operatorname{char}F=p>0$, если не оговорено противное; $\Om$~---
фиксированное алгебраическое замыкание $F$.

\begin{center}
\begin{tikzpicture}[every node/.style={font=\small}]
\begin{scope}
\foreach \i/\lab in {0/{\alpha},1/{\zeta\alpha},2/{\zeta^2\alpha},3/{\zeta^3\alpha},4/{\zeta^4\alpha}}
  \node[circle,fill=cold,inner sep=1.3pt,label={[cold]{90+72*\i}:$\lab$}] (k\i) at ({90+72*\i}:1.15) {};
\foreach \i [evaluate=\i as \j using {int(mod(\i+1,5))}] in {0,...,4}
  \draw[->,cnew,thick,shorten <=3pt,shorten >=3pt] (k\i) to[bend right=22] (k\j);
\node[cnew] at (0,0) {$\sigma$};
\node at (0,-2.05) {Куммер: $\alpha^n=a$, $\sigma\alpha=\zeta\alpha$};
\end{scope}
\begin{scope}[xshift=7cm]
\foreach \i/\lab in {0/{\alpha},1/{\alpha+1},2/{\alpha+2},3/{\alpha+3},4/{\alpha+4}}
  \node[circle,fill=cold,inner sep=1.3pt,label={[cold]{90+72*\i}:$\lab$}] (a\i) at ({90+72*\i}:1.15) {};
\foreach \i [evaluate=\i as \j using {int(mod(\i+1,5))}] in {0,...,4}
  \draw[->,cnew,thick,shorten <=3pt,shorten >=3pt] (a\i) to[bend right=22] (a\j);
\node[cnew] at (0,0) {$\sigma$};
\node at (0,-2.05) {Артин--Шрайер: $\alpha^p-\alpha=a$, $\sigma\alpha=\alpha+1$};
\end{scope}
\end{tikzpicture}
\end{center}
\figcap{Корни неприводимого многочлена $X^5-a$ (при $\zeta=\zeta_5\in F$, $\operatorname{char}F\ne5$) и неприводимого
многочлена $X^5-X-a$ (при $\operatorname{char}F=5$). В обоих случаях образующая $\sigma$ группы Галуа циклически
переставляет корни; в первом случае она умножает корень на $\zeta$, во втором прибавляет к нему $1$.}

\subsection*{Аддитивная теорема Гильберта~90}

\begin{lemma}[5.32a: аддитивная теорема Гильберта~90]
Пусть $E/F$~--- конечное расширение Галуа с группой $G$ (характеристика произвольна).
\begin{enumerate}[label=\textup{(\alph*)}]
\item $H^1(G,E)=0$ для $G$-модуля $(E,+)$ (\rf{5.20}), т.\,е.\ для каждого скрещённого гомоморфизма $f\colon G\to E$
найдётся $\beta\in E$ с $f(\sigma)=\sigma\beta-\beta$ для всех $\sigma\in G$.
\item Если $G=\langle\sigma\rangle$ циклична, то $\operatorname{Tr}_{E/F}(\alpha)=0$ тогда и только тогда, когда
$\alpha=\sigma\beta-\beta$ для некоторого $\beta\in E$.
\end{enumerate}
\end{lemma}

Пункт (a)~--- это теорема \rf{5.23a}, пункт (b)~--- следствие \rf{5.25a} (там $\alpha=\beta'-\sigma\beta'$; положите
$\beta=-\beta'$). Доказательство \rf{5.23a} явное: если $\operatorname{Tr}_{E/F}(\gamma)=1$ (такой $\gamma$ есть по \rf{5.15a}), то
годится $\beta=-\sum_{\tau\in G}f(\tau)\,\tau(\gamma)$. Это аддитивный аналог теоремы \rf{5.23} и следствия \rf{5.25}: там элементы
нормы $1$~--- это элементы вида $\beta/\sigma\beta$, здесь элементы следа $0$~--- элементы вида $\sigma\beta-\beta$. Если
$\operatorname{char}F\nmid n=|G|$, можно взять $\gamma=1/n$ (ср.\ \rf{5.26a}(2)). Нас же будет интересовать случай
$|G|=p=\operatorname{char}F$: тогда $\operatorname{Tr}(1)=p\cdot1=0$, и элемент со следом $1$ приходится искать с помощью
теоремы Дедекинда (\rf{5.15a}).

\subsection*{Многочлены Артина--Шрайера}

\begin{definition}[5.32b]
Для поля $K$ характеристики $p$ положим $\wp(x)=x^p-x$, $x\in K$. Многочлены вида $X^p-X-a$, $a\in K$, называются
\emph{многочленами Артина--Шрайера}.
\end{definition}

Отображение $\wp\colon K\to K$ аддитивно ($(x+y)^p=x^p+y^p$) и даже $\F_p$-линейно ($c^p=c$ для $c\in\F_p$ по малой
теореме Ферма). Многочлен $X^p-X$ имеет $p$ корней $c\in\F_p$, поэтому $X^p-X=\prod_{c\in\F_p}(X-c)$ и
$\ker\wp=\F_p$. Итак, $\wp(K)$~--- подгруппа в $K$, $K/\wp(K)$~--- векторное пространство над $\F_p$, и имеется точная
последовательность
\[
0\longrightarrow\F_p\longrightarrow K\xrightarrow{\ \wp\ }K\longrightarrow K/\wp(K)\longrightarrow0
\]
~--- аналог последовательности $1\to\mu_n\to K^\times\xrightarrow{x\mapsto x^n}K^\times\to K^\times/K^{\times n}\to1$.

\begin{theorem}[5.32c: Артин--Шрайер]
Пусть $\operatorname{char}F=p$, $a\in F$, $f=X^p-X-a$.
\begin{enumerate}[label=\textup{(\alph*)}]
\item Если $\alpha$~--- корень $f$ (в каком-либо расширении), то корни $f$~--- это $p$ различных элементов $\alpha+i$,
$i\in\F_p$; в частности, $f$ сепарабелен.
\item Либо $f$ раскладывается над $F$ на различные линейные множители (это происходит тогда и только тогда, когда
$a\in\wp(F)$), либо $f$ неприводим над $F$. Во втором случае $E=F[\alpha]$~--- поле разложения $f$, расширение $E/F$
циклично степени $p$, и его группа Галуа порождена автоморфизмом $\sigma\colon\alpha\mapsto\alpha+1$.
\item Обратно, если $E/F$~--- циклическое расширение степени $p$ и $\sigma$~--- образующая $\Gal(E/F)$, то
$E=F[\alpha]$, где $\sigma\alpha=\alpha+1$ и $\alpha^p-\alpha=a\in F\smallsetminus\wp(F)$.
\item Пусть $a,b\in F\smallsetminus\wp(F)$ и $\alpha,\beta\in\Om$, $\wp(\alpha)=a$, $\wp(\beta)=b$. Тогда
$F[\alpha]=F[\beta]$ в том и только в том случае, когда $b\in ia+\wp(F)$ для некоторого $i\in\F_p^\times$, т.\,е.\ когда $a$
и $b$ порождают одну и ту же подгруппу в $F/\wp(F)$.
\end{enumerate}
\end{theorem}

\begin{proof}
Пункты (b) и (c) были сформулированы в \rf{5.29}(b); докажем их полностью.

(a) $\wp(\alpha+i)=\wp(\alpha)+\wp(i)=\wp(\alpha)=a$, а элементы $\alpha+i$, $i\in\F_p$, попарно различны. Значит,
$f=\prod_{i\in\F_p}(X-\alpha-i)$. (Иначе: $f'=-1$, и $\gcd(f,f')=1$.)

(b) Пусть $f=gh$ в $F[X]$, где $g$ унитарен и $0<r=\deg g<p$. По (a) $g=\prod_{i\in I}(X-\alpha-i)$ для некоторого
$I\subset\F_p$ из $r$ элементов, и коэффициент $g$ при $X^{r-1}$ равен $-\sum_{i\in I}(\alpha+i)=-r\alpha-\sum_{i\in I}i$.
Он лежит в $F$, а $r\cdot1\ne0$ в $F$ (так как $0<r<p$); значит, $\alpha\in F$, все корни $\alpha+i$ лежат в $F$, и
$a=\wp(\alpha)\in\wp(F)$. Обратно, если $a=\wp(c)$, $c\in F$, то корни $f$~--- это $c+i\in F$. Итак, если $a\notin\wp(F)$,
то $f$ неприводим. Тогда $E=F[\alpha]$ содержит все корни $f$, т.\,е.\ является полем разложения сепарабельного многочлена
$f$; это расширение Галуа степени $p$, и его группа Галуа имеет простой порядок $p$, т.\,е.\ циклична. Элемент $\alpha+1\in
F[\alpha]$~--- корень неприводимого многочлена $f$, поэтому есть $F$-гомоморфизм $\sigma\colon F[\alpha]\to F[\alpha]$,
$\alpha\mapsto\alpha+1$ (\rf{2.1}(b)); он инъективен, $F$-линеен и $[F[\alpha]:F]<\infty$, поэтому $\sigma$~--- автоморфизм;
$\sigma^i\alpha=\alpha+i$, так что $\sigma$ имеет порядок $p$ и порождает группу Галуа.

(c) Так как $\operatorname{Tr}_{E/F}(1)=p\cdot1=0$, по \rf{5.32a}(b) найдётся $\alpha\in E$ с $\sigma\alpha-\alpha=1$. Явно: если
$\operatorname{Tr}_{E/F}(\gamma)=1$, то годится
\[
\alpha=-\sum_{i=0}^{p-1}i\,\sigma^i(\gamma):\qquad
\sigma\alpha=-\sum_{i=0}^{p-1}i\,\sigma^{i+1}(\gamma)=-\sum_{j=0}^{p-1}(j-1)\,\sigma^j(\gamma)=\alpha+\operatorname{Tr}_{E/F}(\gamma)=\alpha+1
\]
(коэффициенты~--- элементы $\F_p$, индексы~--- по модулю $p$, так как $\sigma^p=1$). Поскольку $\sigma\alpha\ne\alpha$, имеем
$\alpha\notin F$, и $F[\alpha]=E$, потому что степень $[E:F]=p$ проста. Далее,
$\sigma(\alpha^p-\alpha)=(\alpha+1)^p-(\alpha+1)=\alpha^p-\alpha$, и $\sigma$ порождает группу Галуа, поэтому
$a=\alpha^p-\alpha\in E^{\langle\sigma\rangle}=F$. Наконец, $a\notin\wp(F)$: иначе по (b) $\alpha\in F$.

(d) Если $b=ia+\wp(c)$, $c\in F$, $i\in\F_p^\times$, то $\wp(i\alpha+c)=i\wp(\alpha)+\wp(c)=b$ ($\wp$ является
$\F_p$-линейным), так что $i\alpha+c$~--- корень $X^p-X-b$, и по (a) $\beta=i\alpha+c+j$ для некоторого $j\in\F_p$. Значит,
$\beta\in F[\alpha]$, и $F[\beta]=F[\alpha]$, так как обе степени равны $p$. Обратно, пусть $F[\alpha]=F[\beta]=E$, и
$\sigma\in\Gal(E/F)$~--- образующая с $\sigma\alpha=\alpha+1$. Элемент $\sigma\beta$~--- корень $X^p-X-b$, поэтому
$\sigma\beta=\beta+i$, $i\in\F_p$, причём $i\ne0$, так как $\beta\notin F$. Тогда
$\sigma(\beta-i\alpha)=\beta+i-i\alpha-i=\beta-i\alpha$, т.\,е.\ $c=\beta-i\alpha\in F$, и
$b=\wp(\beta)=\wp(i\alpha+c)=ia+\wp(c)$.
\end{proof}

Итак, циклические расширения степени $p$ поля $F$ внутри $\Om$ взаимно однозначно соответствуют подгруппам порядка $p$
группы $F/\wp(F)$~--- в точности как в \rf{5.28a} циклические расширения степени $n$ соответствуют циклическим подгруппам
порядка $n$ группы $F^\times/F^{\times n}$. Для $a\in F$ обозначим через $F[\wp^{-1}(a)]$ подполе $\Om$, порождённое корнями
$X^p-X-a$ (по (a) это $F[\alpha]$ для любого из корней $\alpha$).

\subsection*{Абелевы расширения экспоненты $p$}

Расширение Галуа $E/F$ называется абелевым расширением, \emph{экспонента которого делит $p$} (коротко~--- \emph{абелевым
экспоненты $p$}), если $\Gal(E/F)$ абелева и $\sigma^p=1$ для всех $\sigma\in\Gal(E/F)$ (допускается и $E=F$). Конечная абелева группа экспоненты $p$~--- это векторное пространство над
$\F_p$, т.\,е.\ группа вида $(\Z/p\Z)^r$. Для подгруппы $B\subset F$, содержащей $\wp(F)$, положим
$F[\wp^{-1}(B)]=F\bigl[\{\beta\in\Om\mid\wp(\beta)\in B\}\bigr]$.

\begin{theorem}[5.32d]
Отображение $E\mapsto B(E)=F\cap\wp(E)$ задаёт взаимно однозначное соответствие между
\begin{enumerate}[label=\textup{(\alph*)}]
\item конечными абелевыми расширениями $E\subset\Om$ поля $F$ экспоненты $p$ и
\item подгруппами $B\subset F$, содержащими $\wp(F)$, для которых факторгруппа $B/\wp(F)$ конечна.
\end{enumerate}
Обратное отображение: $B\mapsto F[\wp^{-1}(B)]$. При этом $[E:F]=(B(E):\wp(F))$, и спаривание
\[
\Gal(E/F)\times B(E)/\wp(F)\to\F_p,\qquad\langle\sigma,b\rangle=\sigma\beta-\beta\quad(\beta\in E,\ \wp(\beta)=b),
\]
корректно определено, аддитивно по обоим аргументам и совершенно: оно отождествляет каждую из групп с группой
гомоморфизмов другой в $\F_p$.
\end{theorem}

Доказательство повторяет доказательство теоремы Куммера \rf{5.30}; вместо мультипликативной теоремы Гильберта~90
используется аддитивная (\rf{5.32a}), а вместо $\mu_n$~--- ядро $\F_p$ отображения $\wp$.

\begin{proof}
\emph{Шаг 1: для любого конечного расширения Галуа $E/F$ с группой $G$ отображение
\[
\varphi_E\colon B(E)/\wp(F)\to\Hom(G,\F_p),\qquad b\mapsto\bigl(\sigma\mapsto\sigma\beta-\beta\bigr),\quad\wp(\beta)=b,\ \beta\in E,
\]
корректно определено и является изоморфизмом групп.} Пусть $b\in B(E)$, $b=\wp(\beta)$, $\beta\in E$. Элемент $\sigma\beta$~---
тоже корень $X^p-X-b$ (так как $b\in F$), поэтому $\sigma\beta-\beta\in\F_p$ (\rf{5.32c}(a)). Замена $\beta$ на другой корень
$\beta+i$ не меняет разности. Отображение $\sigma\mapsto\sigma\beta-\beta$~--- гомоморфизм:
$\sigma\tau\beta-\beta=\sigma(\tau\beta-\beta)+(\sigma\beta-\beta)=(\tau\beta-\beta)+(\sigma\beta-\beta)$, поскольку
$\tau\beta-\beta\in\F_p$ неподвижен относительно $\sigma$. Если $\wp(\beta)=b$ и $\wp(\beta')=b'$, то
$\wp(\beta+\beta')=b+b'$, так что $\varphi_E$ аддитивно. Ядро: $\sigma\beta=\beta$ для всех $\sigma$ $\iff$
$\beta\in E^G=F$ $\iff$ $b\in\wp(F)$. Сюръективность: гомоморфизм $\chi\colon G\to\F_p\subset E$~--- скрещённый
гомоморфизм ($G$ действует на $\F_p$ тривиально), и по \rf{5.32a}(a) $\chi(\sigma)=\sigma\beta-\beta$ для некоторого
$\beta\in E$. Тогда $\sigma\wp(\beta)=\wp(\beta+\chi(\sigma))=\wp(\beta)$ для всех $\sigma$, т.\,е.\
$b=\wp(\beta)\in F\cap\wp(E)=B(E)$, и $\varphi_E(b)=\chi$. В частности, группа $B(E)/\wp(F)$ конечна.

Если $G$ абелева экспоненты $p$, т.\,е.\ $G\cong\F_p^{\,d}$, то $|\Hom(G,\F_p)|=p^d=|G|$, и шаг 1 даёт $[E:F]=(B(E):\wp(F))$.

\emph{Шаг 2: если $(B:\wp(F))<\infty$, то $E_B=F[\wp^{-1}(B)]$~--- конечное абелево расширение экспоненты $p$.} Пусть
$b_1,\dots,b_r\in B$ порождают $B/\wp(F)$, и $\wp(\beta_j)=b_j$. Каждый $b\in B$ имеет вид $b=\sum n_jb_j+\wp(c)$, $n_j\in\Z$,
$c\in F$; тогда $\beta'=\sum n_j\beta_j+c$ удовлетворяет $\wp(\beta')=b$, и все корни $X^p-X-b$, т.\,е.\ $\beta'+i$, лежат в
$F[\beta_1,\dots,\beta_r]$. Значит, $E_B=F[\beta_1,\dots,\beta_r]$~--- поле разложения многочлена
$\prod_j(X^p-X-b_j)$. Можно считать $b_j$ попарно различными; тогда многочлены $X^p-X-b_j$ не имеют общих корней, и
по \rf{5.32c}(a) их произведение сепарабельно. Поэтому $E_B/F$~--- конечное расширение Галуа (\rf{3.10}). Отображение $\sigma\mapsto(\sigma\beta_j-\beta_j)_{j=1}^r$~--- инъективный гомоморфизм
$\Gal(E_B/F)\to\F_p^{\,r}$ (гомоморфизм~--- по шагу 1; инъективен, так как $\sigma$, оставляющий на месте все $\beta_j$,
тождествен на $E_B$). Значит, $\Gal(E_B/F)$ абелева экспоненты $p$.

\emph{Шаг 3: $E_{B(E)}=E$.} Пусть $E$~--- конечное абелево расширение экспоненты $p$. Ясно, что $E_{B(E)}\subset E$ (если
$X^p-X-b$ имеет корень в $E$, то все его корни в $E$) и $B(E_{B(E)})\supset B(E)$. По шагам 1 и 2
\[
[E:F]\ge[E_{B(E)}:F]=(B(E_{B(E)}):\wp(F))\ge(B(E):\wp(F))=[E:F],
\]
так что всюду равенства, и $E_{B(E)}=E$.

\emph{Шаг 4: $B(E_B)=B$.} Пусть $E=E_B$, $G=\Gal(E/F)$. Ясно, что $B\subset B(E)$. Образ $S=\varphi_E(B/\wp(F))$~---
подпространство в $\Hom(G,\F_p)$. Элемент $\sigma\in G$, на котором обращаются в нуль все $\chi\in S$, оставляет на месте
все $\beta$ с $\wp(\beta)\in B$, а они порождают $E$; значит, $\sigma=1$. Но $G$~--- векторное пространство над $\F_p$
некоторой размерности $d$, $\Hom(G,\F_p)$~--- двойственное пространство, и для подпространства $S$ двойственного
пространства размерность его аннулятора $\{\sigma\mid\chi(\sigma)=0\ \forall\chi\in S\}$ равна $d-\dim S$. Значит, $\dim S=d$,
$S=\Hom(G,\F_p)=\varphi_E(B(E)/\wp(F))$, и из инъективности $\varphi_E$ получаем $B=B(E)$.

Шаги 3 и 4 показывают, что отображения $E\mapsto B(E)$ и $B\mapsto F[\wp^{-1}(B)]$ взаимно обратны. Спаривание
корректно и биаддитивно по шагу 1; отображение $B(E)/\wp(F)\to\Hom(G,\F_p)$ является изоморфизмом по шагу 1, а
отображение $G\to\Hom(B(E)/\wp(F),\F_p)$ инъективно (как в шаге 4) и потому биективно, так как обе группы имеют порядок
$[E:F]$.
\end{proof}

\begin{example}[5.32e: дискриминант Берлекэмпа]
Пусть $\operatorname{char}F=2$ и $f\in F[X]$~--- сепарабельный многочлен степени $n\ge2$ с корнями $\alpha_1,\dots,\alpha_n$
в поле разложения $E$. Обычный дискриминант здесь бесполезен для вопроса <<лежит ли $G_f$ в $A_n$?>>: в характеристике $2$
имеем $\prod_{i<j}(\alpha_i-\alpha_j)=\prod_{i<j}(\alpha_i+\alpha_j)$, а это выражение симметрично по корням, т.\,е.\ лежит в
$F$, так что дискриминант всегда квадрат (\rf{4.4}). Вместо квадратного корня из дискриминанта возьмём
\[
\beta=\sum_{i<j}\frac{\alpha_i}{\alpha_i+\alpha_j}.
\]
Если $\sigma\in G_f$ переставляет корни по перестановке $\pi$, то слагаемое для пары $i<j$ переходит в
$\alpha_{\pi i}/(\alpha_{\pi i}+\alpha_{\pi j})$. При $\pi i<\pi j$ это слагаемое $\beta$ для пары $\{\pi i,\pi j\}$; при
$\pi i>\pi j$ оно отличается от слагаемого для этой пары на $1$, так как
$\frac{\alpha_k}{\alpha_k+\alpha_l}+\frac{\alpha_l}{\alpha_k+\alpha_l}=1$. Поэтому $\sigma\beta=\beta+N(\pi)$, где $N(\pi)$~---
число инверсий $\pi$ (по модулю $2$), т.\,е.\ $\sigma\beta=\beta$ для чётных и $\sigma\beta=\beta+1$ для нечётных $\pi$.
Далее, $\wp(\beta)=\beta^2+\beta=\sum_{i<j}\frac{\alpha_i^2+\alpha_i(\alpha_i+\alpha_j)}{(\alpha_i+\alpha_j)^2}
=\sum_{i<j}\frac{\alpha_i\alpha_j}{(\alpha_i+\alpha_j)^2}=:D$, и $\sigma D=\wp(\sigma\beta)=\wp(\beta)=D$, т.\,е.\ $D\in F$.
Так как $(\alpha_i+\alpha_j)^2=\alpha_i^2+\alpha_j^2$, то $D=\sum_{i<j}\alpha_i\alpha_j/(\alpha_i^2+\alpha_j^2)$~--- это обычная
запись \emph{дискриминанта Берлекэмпа} (ср.\ \rf{4.4}). Итак, $\beta$~--- корень многочлена Артина--Шрайера $X^2+X+D$, и
\[
G_f\subset A_n\iff\beta\in F\iff D\in\wp(F)\iff X^2+X+D\ \text{имеет корень в}\ F.
\]
Например, над $\F_2$ (где $\wp(\F_2)=0$) многочлен $(X+1)(X^2+X+1)$ имеет $G_f=\langle(2\,3)\rangle\not\subset A_3$, и
действительно $D=1$; для неприводимого $X^3+X+1$ группа $G_f\cong\Z/3\subset A_3$ и $D=0$.
\end{example}

\subsection*{Примеры}

\begin{example}[5.32f: конечные поля]
Пусть $F=\F_q$, $q=p^f$, и $\operatorname{Tr}=\operatorname{Tr}_{\F_q/\F_p}$, $\operatorname{Tr}(x)=x+x^p+\dots+x^{p^{f-1}}$ (группа
$\Gal(\F_q/\F_p)$ порождена $x\mapsto x^p$, \rf{4.20}).
\begin{enumerate}[label=\textup{(\alph*)}]
\item \emph{$\wp(\F_q)=\ker\operatorname{Tr}$, и это подгруппа индекса $p$.} Так как $\ker\wp=\F_p$, то $|\wp(\F_q)|=q/p$.
Далее, $\operatorname{Tr}(x^p)=\operatorname{Tr}(x)$ (слагаемые переставляются, так как $x^{p^f}=x$), поэтому
$\operatorname{Tr}(\wp(x))=0$ и $\wp(\F_q)\subset\ker\operatorname{Tr}$. А $\ker\operatorname{Tr}$ состоит из корней многочлена
$X+X^p+\dots+X^{p^{f-1}}$ степени $p^{f-1}$, так что $|\ker\operatorname{Tr}|\le q/p$. Значит, $\wp(\F_q)=\ker\operatorname{Tr}$, и
$\operatorname{Tr}$ индуцирует изоморфизм $\F_q/\wp(\F_q)\cong\F_p$.
\item \emph{Действие Фробениуса.} Пусть $\alpha\in\Om$, $\alpha^p=\alpha+a$, $a\in\F_q$. Возводя в степень $p$, получаем
$\alpha^{p^2}=\alpha+a+a^p$, \dots, $\alpha^{p^k}=\alpha+a+a^p+\dots+a^{p^{k-1}}$, и при $k=f$
\[
\alpha^q=\alpha+\operatorname{Tr}(a).
\]
Значит, $\alpha\in\F_q$ тогда и только тогда, когда $\operatorname{Tr}(a)=0$: \emph{многочлен $X^p-X-a$ неприводим над
$\F_q$ тогда и только тогда, когда $\operatorname{Tr}_{\F_q/\F_p}(a)\ne0$} (в согласии с (a) и \rf{5.32c}(b)). В этом случае
образующая $x\mapsto x^q$ группы $\Gal(\F_{q^p}/\F_q)$ сдвигает корень на $\operatorname{Tr}(a)$, т.\,е.\
$\langle x\mapsto x^q,\,a\rangle=\operatorname{Tr}_{\F_q/\F_p}(a)$ в спаривании \rf{5.32d}.
\item По (a) и \rf{5.32d} у $\F_q$ ровно одно абелево расширение экспоненты $p$ степени $p$ в $\Om$; это $\F_{q^p}$, и оно
получается присоединением корня любого $X^p-X-a$ с $\operatorname{Tr}(a)\ne0$. Это согласуется с \rf{4.23}: в $\Om$ есть ровно одно
поле из $q^p$ элементов.
\end{enumerate}
\end{example}

\begin{example}[5.32g: многочлен $X^p-X-1$ над $\F_p$]
При $q=p$ след тождествен, так что $X^p-X-a$ неприводим над $\F_p$ при всех $a\ne0$, в частности $X^p-X-1$. Пусть
$\alpha$~--- его корень. По \rf{5.32f}(b) $\alpha^{p^i}=\alpha+i$, т.\,е.\ автоморфизм Фробениуса $x\mapsto x^p$ действует как
$\alpha\mapsto\alpha+1$. Перемножим все сопряжённые элементы:
\[
\alpha^{1+p+\dots+p^{p-1}}=\prod_{i=0}^{p-1}\alpha^{p^i}=\prod_{i\in\F_p}(\alpha+i)=\prod_{j\in\F_p}(\alpha-j)=\alpha^p-\alpha=1
\]
(использовано тождество $\prod_{j\in\F_p}(X-j)=X^p-X$). Итак, $\alpha^N=1$ для $N=\frac{p^p-1}{p-1}$; заметим, что
$N\equiv1\pmod p$. Значит, \emph{корни $X^p-X-1$ над $\F_p$~--- корни из единицы степени $N$}, и $X^p-X-1$ делит $X^N-1$.
\begin{itemize}
\item $p=2$: $N=3$, $X^2+X+1=\Phi_3$~--- корни являются первообразными кубическими корнями из единицы.
\item $p=3$: $N=13$, и над $\F_3$ круговой многочлен $\Phi_{13}$ раскладывается в произведение четырёх неприводимых
многочленов степени $3$:
$\Phi_{13}=(X^3-X-1)(X^3+X^2-1)(X^3+X^2+X-1)(X^3-X^2-X-1)$.
\end{itemize}
Порядок $\alpha$ в $\F_{p^p}^\times$ равен в точности $N$ при $p=2,3,5,7$ (при $p=5$: $N=781=11\cdot71$; при $p=7$:
$N=137257=29\cdot4733$). В общем случае из сказанного следует лишь, что порядок делит $N$; для проверки равенства
достаточно убедиться, что $\alpha^{N/\ell}\ne1$ для простых $\ell\mid N$. Кроме того, любое конечное расширение
$\F_{q^m}/\F_q$ порождено образующей циклической группы $\F_{q^m}^\times$, т.\,е.\ корнем из единицы степени $q^m-1$,
взаимно простой с $p$. Поэтому над конечными полями расширения Артина--Шрайера получаются и присоединением корней из
единицы.
\end{example}

\begin{example}[5.32h: поле рациональных функций]
Пусть $k$~--- любое поле характеристики $p$ и $F=k(t)$, $t$ трансцендентно над $k$. Покажем: \emph{если
$a=\sum_{m\in M}c_mt^m\ne0$~--- многочлен от $t$, все показатели которого $m\ge1$ и не делятся на $p$, то
$a\notin\wp(F)$.} Пусть $a=r^p-r$, $r=g/h$, $g,h\in k[t]$ взаимно просты, $h$ унитарен. Тогда $g^p-gh^{p-1}=ah^p$. Если
$h\ne1$, то неприводимый делитель $\pi$ многочлена $h$ делит $g^p$, а значит, и $g$~--- противоречие. Итак, $r=g$~---
многочлен. Если $d=\deg g\ge1$, то коэффициент многочлена $g^p-g$ при $t^{pd}$ равен $g_d^{\,p}\ne0$ ($g_d$~--- старший
коэффициент $g$, а $\deg g<pd$), но у $a$ нет одночленов $t^{pd}$. Если же $g\in k$, то $a=g^p-g$~--- константа, а у $a$
нет свободного члена, т.\,е.\ $a=0$. Противоречие.

Следствия.
\begin{enumerate}[label=\textup{(\alph*)}]
\item Многочлен $X^p-X-t$ неприводим над $k(t)$ для любого поля $k$ характеристики $p$, и его группа Галуа~--- $\Z/p\Z$.
\item Классы $t^m$, $m\ge1$, $p\nmid m$, линейно независимы в $F/\wp(F)$ над $\F_p$ (ненулевая $\F_p$-линейная комбинация
имеет вид из утверждения выше). Поэтому, по \rf{5.32c}(d), поля $F[\alpha_m]$, $\wp(\alpha_m)=t^m$, при таких $m$ попарно
различны, а $F/\wp(F)$~--- бесконечномерное пространство над $\F_p$; в частности, уже $F=\F_p(t)$ (с конечным полем
констант) имеет бесконечно много различных циклических расширений степени $p$. Композит $r$ из них имеет группу Галуа
$(\Z/p\Z)^r$ (\rf{5.32d}).
\item Если поле $k$ совершенно (например, $k=\F_p$ или $\Fpb$), одночлены $ct^{pm}$, $m\ge1$, ничего нового не дают:
$\wp(ct^m)=c^pt^{pm}-ct^m$, т.\,е.\ $c^pt^{pm}\equiv ct^m\pmod{\wp(F)}$, а каждый элемент совершенного $k$ имеет вид $c^p$;
в частности, $t^{pm}\equiv t^m$ (это верно при любом $k$). Например, при $p=3$ поля $F[\wp^{-1}(t)]$ и
$F[\wp^{-1}(t^3)]$ совпадают, а $F[\wp^{-1}(t)]$ и $F[\wp^{-1}(t^2)]$ различны.

Для несовершенного $k$ первое утверждение неверно: при $k=\F_p(s)$ класс $st^p$ не сводится к классам $ct^m$, $p\nmid m$,
и постоянных. Действительно, пусть $st^p-\sum_{p\nmid m}c_mt^m-c_0=g^p-g$, где $c_m,c_0\in k$ и $g\in F$. Как в
доказательстве выше, $g\in k[t]$, и $\deg g\le1$: при $d=\deg g\ge2$ коэффициент правой части при $t^{pd}$ равен
$g_d^{\,p}\ne0$, а у левой части он нулевой. Сравнение коэффициентов при $t^p$ даёт $s=g_1^{\,p}\in k^p$~---
противоречие. Постоянные из $k\smallsetminus\wp(k)$ тоже дают новые расширения, например $X^p-X-1$ над $\F_p(t)$.
\end{enumerate}
\end{example}

\subsection*{Почему одних радикалов не хватает}

Напомним (см.\ определение перед \rf{3.28} и \rf{5.34}), что башня полей $F=F_0\subset F_1\subset\dots\subset F_m$ называется
\emph{радикальной}, если $F_i=F_{i-1}[\alpha_i]$ и $\alpha_i^{n_i}\in F_{i-1}$ для некоторых целых $n_i\ge1$. В
характеристике $p$ шаги с $p\mid n_i$ разрешены, но, как мы видели, корни $p$-й степени дают лишь чисто несепарабельные
расширения.

Видно и то, где ломается доказательство \rf{5.27}. Пусть $E/F$~--- циклическое расширение степени $p$ с образующей $\sigma$.
Резольвента Лагранжа $\sum_i\zeta^i\sigma^i(\gamma)$ строится по первообразному корню $\zeta$ степени $p$ из единицы, а в
характеристике $p$ корень степени $p$ из единицы только один~--- $1$; при $\zeta=1$ резольвента вырождается в след
$\sum_i\sigma^i(\gamma)=\operatorname{Tr}_{E/F}(\gamma)\in F$. Её роль играет элемент $\alpha=-\sum_ii\,\sigma^i(\gamma)$
из \rf{5.32c}(c) (где $\operatorname{Tr}_{E/F}(\gamma)=1$): он не умножается на корень из единицы, а сдвигается,
$\sigma\alpha=\alpha+1$.

Покажем, что циклическое расширение степени $p$ в радикальную башню, вообще говоря, не помещается.

\begin{proposition}[5.32i]
Пусть $\operatorname{char}F=p$, $F=F_0\subset F_1\subset\dots\subset F_m\subset\Om$~--- радикальная башня, $\alpha_i^{n_i}\in F_{i-1}$,
и $F$ содержит первообразный корень степени $\ell$ из единицы для каждого простого $\ell\ne p$, делящего $n_1\cdots n_m$.
Если $F\subset L\subset F_m$ и $L/F$~--- расширение Галуа, то $[L:F]$~--- произведение простых чисел, отличных от $p$. В
частности, если $a\notin\wp(F)$, то корни многочлена $X^p-X-a$ не лежат в $F_m$.
\end{proposition}

\begin{proof}
\emph{Шаг 0: можно считать все $n_i$ простыми.} Если $n_i=1$, то $F_i=F_{i-1}$. Если $n_i=\ell_1\cdots\ell_s$ (простые
множители с повторениями), вставим между $F_{i-1}$ и $F_i$ поля
$F_{i-1}[\alpha_i^{n_i/\ell_1}]\subset F_{i-1}[\alpha_i^{n_i/(\ell_1\ell_2)}]\subset\dots\subset F_{i-1}[\alpha_i]$; на $j$-м шаге
присоединяется элемент, $\ell_j$-я степень которого равна $\alpha_i^{n_i/(\ell_1\cdots\ell_{j-1})}$ и лежит в предыдущем поле.
Множество простых делителей показателей при этом не меняется.

\emph{Шаг 1: два типа шагов.} Пусть $n_i=p$. Для $\gamma=\sum_jc_j\alpha_i^j\in F_i$ ($c_j\in F_{i-1}$) имеем
$\gamma^p=\sum_jc_j^p\alpha_i^{jp}\in F_{i-1}$, поэтому минимальный многочлен $\gamma$ над $F_{i-1}$ делит
$X^p-\gamma^p=(X-\gamma)^p$. Если $\gamma$ сепарабелен над $F_{i-1}$, этот минимальный многочлен равен $X-\gamma$, т.\,е.\
$\gamma\in F_{i-1}$: \emph{каждый элемент $F_i$, сепарабельный над $F_{i-1}$, лежит в $F_{i-1}$.} Пусть теперь $n_i=\ell\ne p$
и $\zeta\in F$~--- первообразный корень степени $\ell$ из единицы. Если $b=\alpha_i^\ell=0$, то $F_i=F_{i-1}$. Иначе корни
$X^\ell-b$~--- это $\ell$ различных элементов $\zeta^j\alpha_i\in F_i$, поэтому $F_i$~--- поле разложения сепарабельного
многочлена $X^\ell-b$ над $F_{i-1}$, а $\sigma\mapsto\sigma\alpha_i/\alpha_i$~--- инъективный гомоморфизм
$\Gal(F_i/F_{i-1})\to\mu_\ell$ (как перед \rf{5.27}). Значит, $[F_i:F_{i-1}]\in\{1,\ell\}$.

\emph{Шаг 2: подсчёт степени.} Поле $L$~--- поле разложения над $F$ некоторого сепарабельного многочлена $f$. Положим
$L_i=LF_i\subset F_m$ (композит) и $d_i=[L_i:F_i]$; поле $L_i$~--- поле разложения $f$ над $F_i$, так что $L_i/F_i$~---
расширение Галуа, $d_0=[L:F]$ и $d_m=1$ (так как $L\subset F_m$). Применим \rf{3.19} к расширению Галуа $L_{i-1}/F_{i-1}$ и
расширению $F_i/F_{i-1}$: композит $L_{i-1}F_i$ равен $L_i$, и $[L_i:F_i]=[L_{i-1}:L_{i-1}\cap F_i]$. Поэтому
\[
d_{i-1}=[L_{i-1}:L_{i-1}\cap F_i]\,[L_{i-1}\cap F_i:F_{i-1}]=d_i\,e_i,\qquad e_i=[L_{i-1}\cap F_i:F_{i-1}].
\]
Поле $L_{i-1}\cap F_i$ сепарабельно над $F_{i-1}$ (оно лежит в расширении Галуа $L_{i-1}$), а $e_i$ делит $[F_i:F_{i-1}]$.
По шагу 1 $e_i=1$ при $n_i=p$ и $e_i\in\{1,\ell\}$ при $n_i=\ell\ne p$. Итак, $[L:F]=d_0=e_1e_2\cdots e_m$~--- произведение
простых, отличных от $p$.

Наконец, если $\wp(\alpha)=a\notin\wp(F)$ и $\alpha\in F_m$, то $L=F[\alpha]\subset F_m$~--- расширение Галуа степени $p$
(\rf{5.32c}), что противоречит доказанному.
\end{proof}

\begin{corollary}[5.32j]
Пусть $k$~--- любое поле характеристики $p$ (например, $\F_p$ или $\Fpb$), $F=k(t)$, и $\alpha\in\Om$, $\alpha^p-\alpha=t$.
Тогда $F[\alpha]/F$~--- циклическое расширение степени $p$, но $\alpha$ не лежит ни в какой радикальной башне над $F$:
уравнение $X^p-X-t=0$ с группой Галуа $\Z/p\Z$ не разрешимо в радикалах.
\end{corollary}

\begin{proof}
Цикличность~--- \rf{5.32h}(a) и \rf{5.32c}(b). Пусть $\alpha\in F_m$ для радикальной башни $F_i=F_{i-1}[\alpha_i]$. Пусть $N$~---
произведение всех простых $\ell\ne p$, делящих $n_1\cdots n_m$ ($N=1$, если таких нет), и $\zeta\in\Om$~--- первообразный
корень степени $N$ из единицы (он существует, так как $p\nmid N$, \rf{5.8}). Поля $F'_i=F[\zeta,\alpha_1,\dots,\alpha_i]$ образуют
радикальную башню над $F'=F[\zeta]$, $F'$ содержит первообразные корни $\zeta^{N/\ell}$ степени $\ell$ из единицы, и
$\alpha\in F'_m$. Но $F'=k'(t)$, где $k'=k[\zeta]$, и $t$ трансцендентно над $k'$ (поле $k'$ алгебраично над $k$). По
\rf{5.32h} $t\notin\wp(k'(t))$, и \rf{5.32i} даёт противоречие.
\end{proof}

Доказательство \rf{5.32j} даёт больше. Пусть $F$~--- любое поле характеристики $p$, $a\in F\smallsetminus\wp(F)$, и корень
$\alpha$ многочлена $X^p-X-a$ лежит в радикальной башне $F=F_0\subset\dots\subset F_m$, $\alpha_i^{n_i}\in F_{i-1}$; пусть
$N$~--- произведение простых $\ell\ne p$, делящих $n_1\cdots n_m$ ($N=1$, если таких нет), и $\zeta_N\in\Om$~---
первообразный корень степени $N$ из единицы. Тогда $\alpha\in F[\zeta_N]$. Иначе ни один из корней $\alpha+i$ не лежит в $F'=F[\zeta_N]$, т.\,е.\
$a\notin\wp(F')$, и \rf{5.32i}, применённое (как в доказательстве \rf{5.32j}) к радикальной башне $F'_i=F'[\alpha_1,\dots,\alpha_i]$
над $F'$, показывает, что $\alpha\notin F'_m$,~--- противоречие, так как $\alpha\in F_m\subset F'_m$. Обратно,
$F\subset F[\zeta_N]$~--- радикальная башня из одного шага ($\zeta_N^N=1\in F$), так что элементы $F[\zeta_N]$ выражаются в
радикалах. Итак, \emph{корни многочлена Артина--Шрайера выражаются в радикалах \textup(лежат в некоторой радикальной
башне над $F$\textup) тогда и только тогда, когда они лежат в круговом расширении $F[\zeta_N]$ для некоторого $N$, не
делящегося на $p$}.

Например, над $\F_p$ корень $X^p-X-1$~--- корень из единицы (\rf{5.32g}), т.\,е.\ радикал; то же верно над $\F_p(t)$ для
$X^p-X-c$, $c\in\F_p^\times$: его корни лежат в $\F_{p^p}(t)$, а поле $\F_{p^p}$ порождено над $\F_p$ корнем из единицы.
А корни $X^p-X-t$ не лежат ни в каком поле $k(t)[\zeta_N]=k[\zeta_N](t)$ (\rf{5.32h}, как в доказательстве \rf{5.32j}):
препятствие связано именно с тем, что $t$ <<не постоянна>>. Поэтому в характеристике $p$ к радикалам приходится добавлять
корни многочленов Артина--Шрайера; с этой поправкой теорема Галуа о разрешимости остаётся верной (\rf{5.35b}).

\subsection*{Векторы Витта}

Циклическое расширение степени $p^2$ в характеристике $p$ строится двумя шагами Артина--Шрайера, но не всякая такая
башня циклична: например, $\F_p(t)[\wp^{-1}(t),\wp^{-1}(t^2)]$ при $p\ne2$ имеет группу $(\Z/p\Z)^2$ (\rf{5.32h}). Чтобы описать
циклические расширения степени $p^n$, нужно заменить поле $F$ и группу $\F_p$ кольцом $W_n(F)$ и группой
$W_n(\F_p)\cong\Z/p^n\Z$. Образец~--- кольцо $\Zp=\varlim\Z/p^n\Z$ (\rf{7.33}): его элементы задаются последовательностями
<<цифр>> из $\F_p$, но сложение цифр идёт с переносом. Векторы Витта делают <<перенос>> многочленом с целыми
коэффициентами, пригодным для любого коммутативного кольца.

Фиксируем простое $p$. \emph{Многочлены Витта}~--- это
\[
w_n=\sum_{i=0}^np^iX_i^{p^{n-i}}=X_0^{p^n}+pX_1^{p^{n-1}}+\dots+p^nX_n\in\Z[X_0,\dots,X_n],\qquad n\ge0;
\]
$w_0=X_0$, $w_1=X_0^p+pX_1$, $w_2=X_0^{p^2}+pX_1^p+p^2X_2$. Для коммутативного кольца $A$ и последовательности
$a=(a_0,a_1,\dots)$ элементов $A$ значения $w_n(a)=w_n(a_0,\dots,a_n)$ называются \emph{призрачными компонентами} $a$.
Если $p$ обратимо в $A$, то по $(w_0(a),w_1(a),\dots)$ последовательность $a$ восстанавливается рекуррентно:
$a_n=p^{-n}\bigl(w_n(a)-\sum_{i<n}p^ia_i^{p^{n-i}}\bigr)$; если в $A$ нет $p$-кручения ($px=0\Rightarrow x=0$), то $a$
определяется своими призрачными компонентами однозначно. Мы хотим ввести на последовательностях сложение и умножение,
которые в призрачных компонентах становятся покомпонентными. Над $\Q$ формулы однозначно определены; главное~--- что их
коэффициенты целые. Это вытекает из следующей леммы.

\begin{lemma}[5.32k]
Пусть $A$~--- коммутативное кольцо.
\begin{enumerate}[label=\textup{(\alph*)}]
\item Если $a\equiv b\pmod{p^kA}$, $k\ge1$, то $a^p\equiv b^p\pmod{p^{k+1}A}$; следовательно,
$a^{p^m}\equiv b^{p^m}\pmod{p^{k+m}A}$ для всех $m\ge0$.
\item Пусть $\varphi\colon A\to A$~--- эндоморфизм колец, для которого $\varphi(x)\equiv x^p\pmod{pA}$ при всех $x\in A$.
Если $x_n=w_n(a_0,\dots,a_n)$ для некоторых $a_i\in A$, то $x_n\equiv\varphi(x_{n-1})\pmod{p^nA}$ при всех $n\ge1$. Если в
$A$ нет $p$-кручения, верно и обратное: для каждой последовательности $(x_n)_{n\ge0}$ с $x_n\equiv\varphi(x_{n-1})\pmod{p^nA}$
($n\ge1$) существует единственная последовательность $(a_n)$ элементов $A$ с $w_n(a)=x_n$ для всех $n$.
\end{enumerate}
\end{lemma}

\begin{proof}
(a) Пусть $a=b+p^kc$. Тогда $a^p=b^p+\sum_{j=1}^p\binom pj b^{p-j}p^{jk}c^j$; слагаемое с $j=1$ делится на $p\cdot p^k$,
а при $j\ge2$ имеем $jk\ge2k\ge k+1$. Второе утверждение~--- индукция по $m$.

(b) Пусть $i<n$. Из $\varphi(a_i)\equiv a_i^p\pmod p$ по (a) (с $k=1$, $m=n-1-i$) получаем
$\varphi(a_i)^{p^{n-1-i}}\equiv a_i^{p^{n-i}}\pmod{p^{n-i}}$, т.\,е.
\[
p^i\varphi(a_i)^{p^{n-1-i}}\equiv p^ia_i^{p^{n-i}}\pmod{p^n}.\tag{$*$}
\]
Так как $\varphi$~--- гомоморфизм колец, $\varphi(w_{n-1}(a))=\sum_{i<n}p^i\varphi(a_i)^{p^{n-1-i}}$, и по $(*)$
\[
\varphi\bigl(w_{n-1}(a)\bigr)\equiv\sum_{i<n}p^ia_i^{p^{n-i}}=w_n(a)-p^na_n\equiv w_n(a)\pmod{p^n}.
\]
Это первое утверждение. Обратно, пусть $A$ без $p$-кручения и $x_n\equiv\varphi(x_{n-1})\pmod{p^n}$. Положим $a_0=x_0$ и
предположим, что $a_0,\dots,a_{n-1}$ уже найдены, $w_i(a)=x_i$ при $i<n$. По $(*)$
\[
x_n-\sum_{i<n}p^ia_i^{p^{n-i}}\equiv x_n-\varphi(x_{n-1})\equiv0\pmod{p^n},
\]
т.\,е.\ эта разность равна $p^na_n$ для некоторого $a_n\in A$, и тогда $w_n(a)=x_n$. Единственность $a_n$ следует из
отсутствия $p$-кручения.
\end{proof}

\begin{theorem}[5.32l: Витт]
Существуют единственные последовательности многочленов $S_n,P_n\in\Z[X_0,\dots,X_n,Y_0,\dots,Y_n]$ и
$I_n\in\Z[X_0,\dots,X_n]$, $n\ge0$, такие что для всех $n$
\[
w_n(S_0,\dots,S_n)=w_n(X)+w_n(Y),\quad w_n(P_0,\dots,P_n)=w_n(X)\,w_n(Y),\quad w_n(I_0,\dots,I_n)=-w_n(X).
\]
При этом $S_n=X_n+Y_n+C_n$, где $C_n$~--- многочлен от $X_0,\dots,X_{n-1},Y_0,\dots,Y_{n-1}$.
\end{theorem}

\begin{proof}
В кольце $\Q[X,Y]$ (многочлены от всех $X_i,Y_i$) равенство для $S_n$ однозначно определяет $S_n$ по $S_0,\dots,S_{n-1}$:
\[
S_n=p^{-n}\Bigl(w_n(X)+w_n(Y)-\sum_{i<n}p^iS_i^{p^{n-i}}\Bigr),
\]
и индукцией по $n$ видно, что $S_n$ зависит лишь от переменных с индексами $\le n$ и что $S_n-X_n-Y_n$ зависит лишь от
переменных с индексами $<n$ (вклад $p^nX_n+p^nY_n$ в $w_n(X)+w_n(Y)$~--- единственный, содержащий $X_n,Y_n$). То же для
$P_n$ и $I_n$. Остаётся доказать целочисленность.

Пусть $A=\Z[X_0,X_1,\dots,Y_0,Y_1,\dots]$ и $\varphi$~--- эндоморфизм $A$ с $\varphi(X_i)=X_i^p$, $\varphi(Y_i)=Y_i^p$. Для
$g\in A$ имеем $\varphi(g)\equiv g^p\pmod{pA}$: в кольце $A/pA=\F_p[X,Y]$ возведение в степень $p$~--- гомоморфизм, а
$c^p=c$ для $c\in\F_p$, так что $g^p\equiv\sum_\mu c_\mu\mu^p=\varphi(g)$, где $g=\sum c_\mu\mu$~--- разложение по одночленам.
В $A$ нет $p$-кручения. По лемме \rf{5.32k}(b), применённой к последовательностям $X=(X_i)$ и $Y=(Y_i)$,
$w_n(X)\equiv\varphi(w_{n-1}(X))$ и $w_n(Y)\equiv\varphi(w_{n-1}(Y))\pmod{p^n}$. Поэтому последовательность
$x_n=w_n(X)+w_n(Y)$ тоже удовлетворяет $x_n\equiv\varphi(x_{n-1})\pmod{p^n}$, а значит, $x_n=w_n(s)$ для некоторых
$s_n\in A$. В $\Q[X,Y]$ решение единственно, поэтому $s_n=S_n$, т.\,е.\ $S_n\in A$. Для произведения:
$w_n(X)w_n(Y)\equiv\varphi(w_{n-1}(X))\varphi(w_{n-1}(Y))=\varphi\bigl(w_{n-1}(X)w_{n-1}(Y)\bigr)\pmod{p^n}$, для
противоположного элемента~--- $-w_n(X)\equiv\varphi(-w_{n-1}(X))$. Остальное так же.
\end{proof}

Первые многочлены легко найти. Так как $w_0=X_0$, то $S_0=X_0+Y_0$ и $P_0=X_0Y_0$. Из
$S_0^p+pS_1=X_0^p+pX_1+Y_0^p+pY_1$ получаем $S_1=X_1+Y_1+\frac{X_0^p+Y_0^p-(X_0+Y_0)^p}{p}$, а из
$P_0^p+pP_1=(X_0^p+pX_1)(Y_0^p+pY_1)$~--- формулу для $P_1$:
\[
S_0=X_0+Y_0,\quad S_1=X_1+Y_1-\sum_{j=1}^{p-1}\frac1p\binom pjX_0^jY_0^{p-j},\qquad P_0=X_0Y_0,\quad
P_1=X_0^pY_1+X_1Y_0^p+pX_1Y_1.
\]
Например, при $p=2$: $S_1=X_1+Y_1-X_0Y_0$; при $p=3$: $S_1=X_1+Y_1-X_0^2Y_0-X_0Y_0^2$. Для нечётного $p$ имеем
$I_n=-X_n$ (так как $w_n(-X)=-w_n(X)$), а для $p=2$: $I_0=-X_0$, $I_1=-X_1-X_0^2$, $I_2=-X_2-X_1^2-X_0^2X_1-X_0^4$.

\begin{definition}[5.32m]
Пусть $A$~--- коммутативное кольцо. \emph{Кольцо векторов Витта} $W(A)$~--- множество последовательностей
$a=(a_0,a_1,a_2,\dots)$ элементов $A$ с операциями
\[
a+b=\bigl(S_0(a,b),S_1(a,b),\dots\bigr),\qquad ab=\bigl(P_0(a,b),P_1(a,b),\dots\bigr).
\]
Так как $S_k,P_k$ зависят лишь от переменных с индексами $\le k$, те же формулы задают операции на векторах длины $n$
$(a_0,\dots,a_{n-1})$; полученное множество обозначается $W_n(A)$ (\emph{векторы Витта длины $n$}).
\end{definition}

\begin{proposition}[5.32n]
\begin{enumerate}[label=\textup{(\alph*)}]
\item $W(A)$ и $W_n(A)$~--- коммутативные кольца с нулём $(0,0,\dots)$, единицей $(1,0,0,\dots)$ и противоположным
элементом $-a=(I_0(a),I_1(a),\dots)$.
\item Отображение $a\mapsto(w_0(a),w_1(a),\dots)$~--- гомоморфизм колец $W(A)\to A^{\mathbb N}$ (произведение колец); оно
инъективно, если в $A$ нет $p$-кручения, и биективно, если $p$ обратимо в $A$.
\item Гомоморфизм колец $h\colon A\to B$ индуцирует гомоморфизм $W(h)\colon W(A)\to W(B)$, $(a_n)\mapsto(h(a_n))$.
\item Отбрасывание компонент $W(A)\to W_n(A)$, $W_{n+1}(A)\to W_n(A)$~--- сюръективные гомоморфизмы колец; $W_1(A)=A$,
и $W(A)=\varlim_nW_n(A)$.
\end{enumerate}
\end{proposition}

\begin{proof}
По построению $w_n(a+b)=w_n(a)+w_n(b)$ и $w_n(ab)=w_n(a)w_n(b)$ для любого кольца $A$: это тождества между многочленами с
целыми коэффициентами, и в них можно подставить элементы $A$. Пусть $A_0=\Z[X_i,Y_i,Z_i\mid i\ge0]$ и
$X=(X_i),Y=(Y_i),Z=(Z_i)\in W(A_0)$. Каждая аксиома кольца для $W$~--- равенство двух векторов, компоненты которых~---
многочлены с целыми коэффициентами от $X,Y,Z$; например, ассоциативность сложения означает $(X+Y)+Z=X+(Y+Z)$ в $W(A_0)$.
Призрачные компоненты обеих частей одинаковы ($w_n(X)+w_n(Y)+w_n(Z)$), а в $A_0$ нет $p$-кручения, поэтому обе части
равны, т.\,е.\ соответствующие многочлены совпадают. Подставив вместо $X_i,Y_i,Z_i$ элементы любого кольца $A$, получаем
аксиому в $W(A)$. Так же проверяются коммутативность, дистрибутивность, свойства нуля и единицы
($w_n(1,0,0,\dots)=1$ для всех $n$) и противоположного элемента. Это доказывает (a) для $W(A)$ и (b) (инъективность и
биективность~--- см.\ выше). (c): компоненты суммы и произведения~--- многочлены с целыми коэффициентами от компонент
слагаемых. (d): компоненты с номерами $<n$ суммы и произведения зависят только от компонент с номерами $<n$; поэтому
отбрасывание компонент сохраняет операции, и $W_n(A)$~--- кольцо как образ кольца $W(A)$. Равенство $W_1(A)=A$
означает $S_0=X_0+Y_0$, $P_0=X_0Y_0$. Наконец, как множество $W(A)=A^{\mathbb N}$~--- предел множеств $A^n$, а операции
согласованы с отбрасываниями.
\end{proof}

\begin{proposition}[5.32o: Тейхмюллер, сдвиг, Фробениус]
Пусть $A$~--- коммутативное кольцо.
\begin{enumerate}[label=\textup{(\alph*)}]
\item \emph{Представитель Тейхмюллера} $[x]=(x,0,0,\dots)$ мультипликативен: $[xy]=[x][y]$; $w_n([x])=x^{p^n}$.
\item \emph{Сдвиг} $V(a_0,a_1,\dots)=(0,a_0,a_1,\dots)$ аддитивен, и $w_0(Va)=0$, $w_n(Va)=p\,w_{n-1}(a)$ при $n\ge1$. Он
индуцирует инъективный аддитивный гомоморфизм $V\colon W_{n-1}(A)\to W_n(A)$, образ которого~--- ядро гомоморфизма колец
$W_n(A)\to W_1(A)=A$, $a\mapsto a_0$.
\item Существуют многочлены $\mathrm{Fr}_n\in\Z[X_0,\dots,X_{n+1}]$, $n\ge0$, такие что отображение
$\operatorname{Fr}\colon W(A)\to W(A)$, $\operatorname{Fr}(a)=(\mathrm{Fr}_0(a),\mathrm{Fr}_1(a),\dots)$,~--- эндоморфизм
кольца с $w_n(\operatorname{Fr}a)=w_{n+1}(a)$; при этом $\operatorname{Fr}(Va)=p\cdot a$ для всех $a\in W(A)$ (здесь
$p\cdot a=a+\dots+a$).
\item Если $pA=0$ (т.\,е.\ $A$~--- $\F_p$-алгебра), то $\mathrm{Fr}_n\equiv X_n^p\pmod p$, так что
\[
\operatorname{Fr}(a_0,a_1,\dots)=(a_0^p,a_1^p,\dots),\qquad p\cdot a=V\operatorname{Fr}(a)=\operatorname{Fr}V(a)=(0,a_0^p,a_1^p,\dots),
\]
и, по индукции, $p^k\cdot a=(0,\dots,0,a_0^{p^k},a_1^{p^k},\dots)$ ($k$ нулей). Эти формулы верны и в $W_n(A)$.
\end{enumerate}
\end{proposition}

Для произвольного $A$ равенство $V\operatorname{Fr}=p$ неверно: при $A=\Z$ нулевая компонента $p\cdot1$ равна $p$, а у
$V\operatorname{Fr}(1)$ она равна $0$. Верно лишь $\operatorname{Fr}V=p$. Отметим также, что $[x]+[y]\ne[x+y]$ в общем случае.

\begin{proof}
Во всех пунктах используется тот же приём, что и в \rf{5.32n}: тождество проверяется в призрачных компонентах для векторов из
переменных (кольцо многочленов над $\Z$ без $p$-кручения), а затем подставляются элементы $A$.

(a) $w_n([x])=x^{p^n}$, поэтому $w_n([X][Y])=X^{p^n}Y^{p^n}=w_n([XY])$.

(b) $w_n(Va)=\sum_{i=1}^np^ia_{i-1}^{p^{n-i}}=p\sum_{j=0}^{n-1}p^ja_j^{p^{n-1-j}}=p\,w_{n-1}(a)$. Значит,
$w(V(X+Y))=w(VX)+w(VY)$, т.\,е.\ $V(X+Y)=VX+VY$. Компоненты $Va$ с номерами $<n$ зависят лишь от $a_0,\dots,a_{n-2}$, отсюда
$V\colon W_{n-1}(A)\to W_n(A)$; его образ~--- векторы с $a_0=0$.

(c) В $A_0=\Z[X_0,X_1,\dots]$ с эндоморфизмом $\varphi(X_i)=X_i^p$ последовательность $x_n=w_{n+1}(X)$ удовлетворяет
$x_n\equiv\varphi(x_{n-1})\pmod{p^{n+1}}$ (\rf{5.32k}(b) для вектора $X$), и тем более по модулю $p^n$. По \rf{5.32k}(b) найдутся
$\mathrm{Fr}_n\in A_0$ с $w_n(\mathrm{Fr}_0,\dots,\mathrm{Fr}_n)=w_{n+1}(X)$; из рекуррентной формулы видно, что
$\mathrm{Fr}_n$ зависит лишь от $X_0,\dots,X_{n+1}$.
Сдвиг призрачных компонент $(x_0,x_1,\dots)\mapsto(x_1,x_2,\dots)$~--- эндоморфизм кольца $A^{\mathbb N}$, поэтому
$\operatorname{Fr}$~--- эндоморфизм колец. Далее, $w_n(\operatorname{Fr}Va)=w_{n+1}(Va)=p\,w_n(a)=w_n(p\cdot a)$.

(d) Докажем $\mathrm{Fr}_n\equiv X_n^p\pmod{pA_0}$ индукцией по $n$; $\mathrm{Fr}_0=X_0^p+pX_1$. Пусть
$\mathrm{Fr}_i=X_i^p+pG_i$ при $i<n$. По \rf{5.32k}(a) $\mathrm{Fr}_i^{\,p^{n-i}}\equiv X_i^{p^{n+1-i}}\pmod{p^{n-i+1}}$, т.\,е.\
$p^i\mathrm{Fr}_i^{\,p^{n-i}}\equiv p^iX_i^{p^{n+1-i}}\pmod{p^{n+1}}$. Сравним равенство
$\sum_{i=0}^np^i\mathrm{Fr}_i^{\,p^{n-i}}=\sum_{i=0}^{n+1}p^iX_i^{p^{n+1-i}}$ по модулю $p^{n+1}$:
$p^n\mathrm{Fr}_n\equiv p^nX_n^p+p^{n+1}X_{n+1}\equiv p^nX_n^p$. Значит, $p^n(\mathrm{Fr}_n-X_n^p)=p^{n+1}h$ для некоторого
$h\in A_0$, и, так как в $A_0$ нет $p$-кручения, $\mathrm{Fr}_n-X_n^p=ph$. Если $pA=0$, то отсюда $\mathrm{Fr}_n(a)=a_n^p$.
Тогда $p\cdot a=\operatorname{Fr}(Va)=(0,a_0^p,a_1^p,\dots)=V\operatorname{Fr}(a)$.
Так как все операции согласованы с отбрасыванием компонент, формулы верны и в $W_n(A)$.
\end{proof}

\begin{proposition}[5.32p]
$W_n(\F_p)\cong\Z/p^n\Z$ и $W(\F_p)\cong\Zp$ как кольца.
\end{proposition}

\begin{proof}
В $W_n(\F_p)$ ровно $p^n$ элементов. По \rf{5.32o}(d) $p^k\cdot1=(0,\dots,0,1,0,\dots,0)$ (единица на месте с номером $k$), так
как $1^{p^k}=1$. Поэтому $p^{n-1}\cdot1\ne0$, а $p^n\cdot1=0$ в $W_n(\F_p)$, т.\,е.\ аддитивный порядок единицы равен $p^n$.
Гомоморфизм колец $\Z\to W_n(\F_p)$, $m\mapsto m\cdot1$, имеет ядро $p^n\Z$ и, по подсчёту элементов, сюръективен. Эти
изоморфизмы $\Z/p^n\Z\cong W_n(\F_p)$ согласованы с отбрасыванием компонент $W_{n+1}(\F_p)\to W_n(\F_p)$ (это гомоморфизмы
колец), поэтому $W(\F_p)=\varlim W_n(\F_p)\cong\varlim\Z/p^n\Z=\Zp$ (\rf{7.33}).
\end{proof}

Например, в $W_2(\F_2)\cong\Z/4\Z$: $0=(0,0)$, $1=(1,0)$, $2=(0,1)$, $3=(1,1)$ (по формуле $S_1=X_1+Y_1-X_0Y_0$). Так как
$[x]^p=[x^p]=[x]$ для $x\in\F_p$, образ $[x]$ в $\Zp$~--- корень уравнения $y^p=y$, сравнимый с $x$ по модулю $p$; при
$x\ne0$ это тот единственный элемент группы $\mu_{p-1}\subset\Zp$ (см.\ после \rf{7.33a}), который сравним с $x$. Например, в
$W_2(\F_5)\cong\Z/25\Z$: $[2]=7$, $[3]=18$, $[4]=24=-1$, и действительно $7^4=2401\equiv1\pmod{25}$.

\begin{aside}[5.32q]
Аналогично для $q=p^f$ кольцо $W(\F_q)$~--- это кольцо целых неразветвлённого расширения степени $f$ поля $\Qp$ (оно
изоморфно $\Zp[\zeta_{q-1}]$); элементы $\F_q$ при этом поднимаются в $W(\F_q)$ представителями Тейхмюллера. Мы этим
фактом не пользуемся и не доказываем его. Из \rf{5.32o}(d) легко видеть лишь, что в $W(\F_q)$ умножение на $p$ инъективно
($p\cdot a=V\operatorname{Fr}a$, а $V$ и $\operatorname{Fr}$ инъективны: $x\mapsto x^p$ инъективно на поле) и что
$pW(\F_q)=VW(\F_q)$ (здесь нужна совершенность $\F_q$: $\operatorname{Fr}$ сюръективно), откуда
$W(\F_q)/pW(\F_q)=W(\F_q)/VW(\F_q)\cong\F_q$.
\end{aside}

\subsection*{Теория Артина--Шрайера--Витта}

Пусть $K$~--- поле характеристики $p$ и $n\ge1$. По \rf{5.32o}(d) отображение $\operatorname{Fr}(a)=(a_0^p,\dots,a_{n-1}^p)$~---
эндоморфизм кольца $W_n(K)$. Положим
\[
\wp=\operatorname{Fr}-1\colon W_n(K)\to W_n(K),\qquad\wp(a)=\operatorname{Fr}(a)-a
\]
(вычитание~--- в кольце $W_n(K)$). Это гомоморфизм аддитивных групп, и $\ker\wp=\{a\mid a_i^p=a_i\ \forall i\}=W_n(\F_p)\cong\Z/p^n\Z$
(\rf{5.32p}). При $n=1$ получаем прежнее $\wp(x)=x^p-x$ на $W_1(K)=K$. Если $E/F$~--- расширение Галуа с группой $G$, то $G$
действует на $W_n(E)$ покомпонентно автоморфизмами колец (\rf{5.32n}(c)); это действие перестановочно с $\operatorname{Fr}$, $V$
и $\wp$, и $W_n(E)^G=W_n(F)$.

\begin{lemma}[5.32r]
Если $E/F$~--- конечное расширение Галуа с группой $G$, то $H^1(G,W_n(E))=0$.
\end{lemma}

\begin{proof}
Индукция по $n$; при $n=1$ это \rf{5.32a}. Используем точную последовательность $G$-модулей
\[
\begin{tikzcd}
0\arrow[r] & W_{n-1}(E)\arrow[r,"V"] & W_n(E)\arrow[r,"\mathrm{pr}"] & E\arrow[r] & 0,
\end{tikzcd}\qquad\mathrm{pr}(a)=a_0
\]
(\rf{5.32o}(b)). Пусть $f\colon G\to W_n(E)$~--- скрещённый гомоморфизм. Тогда $\mathrm{pr}\circ f$~--- скрещённый гомоморфизм
$G\to E$ ($\mathrm{pr}$ аддитивно и перестановочно с $G$), и по \rf{5.32a} $\mathrm{pr}f(\sigma)=\sigma b_0-b_0$ для некоторого
$b_0\in E$. Положим $b=[b_0]=(b_0,0,\dots,0)$ и $f'(\sigma)=f(\sigma)-(\sigma b-b)$. Это скрещённый гомоморфизм с
$\mathrm{pr}f'(\sigma)=0$, т.\,е.\ $f'(\sigma)=V(g(\sigma))$, $g(\sigma)\in W_{n-1}(E)$. Так как $V$ инъективно, аддитивно и
перестановочно с $G$, $g$~--- скрещённый гомоморфизм, и по предположению индукции $g(\sigma)=\sigma c-c$, $c\in W_{n-1}(E)$.
Тогда $f(\sigma)=\sigma b-b+V(\sigma c-c)=\sigma(b+Vc)-(b+Vc)$.
\end{proof}

\begin{theorem}[5.32s: Артин--Шрайер--Витт]
Пусть $\operatorname{char}F=p$, $n\ge1$, $a=(a_0,\dots,a_{n-1})\in W_n(F)$.
\begin{enumerate}[label=\textup{(\alph*)}]
\item Существует $\beta\in W_n(\Om)$ с $\wp(\beta)=a$; все такие $\beta$ имеют вид $\beta+c$, $c\in W_n(\F_p)$, и поле
$F(\beta)=F(\beta_0,\dots,\beta_{n-1})$ зависит только от $a$. Расширение $F(\beta)/F$~--- конечное расширение Галуа, и
$\sigma\mapsto\sigma\beta-\beta$~--- инъективный гомоморфизм $\Gal(F(\beta)/F)\to W_n(\F_p)\cong\Z/p^n\Z$; в частности,
$\Gal(F(\beta)/F)$ циклична порядка, делящего $p^n$.
\item $[F(\beta):F]=p^n$ тогда и только тогда, когда $a_0\notin\wp(F)$.
\item Каждое циклическое расширение $E\subset\Om$ поля $F$ степени $p^n$ имеет вид $F(\beta)$, $\wp(\beta)=a$, для некоторого
$a\in W_n(F)$ с $a_0\notin\wp(F)$.
\end{enumerate}
\end{theorem}

\begin{proof}
(a) Равенство $\operatorname{Fr}(\beta)=\beta+a$ по \rf{5.32l} означает, что $\beta_k^p=\beta_k+a_k+C_k(\beta_0,\dots,\beta_{k-1},a_0,\dots,a_{k-1})$
для $k=0,\dots,n-1$. Поэтому $\beta$ строится по шагам: $\beta_k$~--- корень многочлена Артина--Шрайера $X^p-X-c_k$ с
$c_k=a_k+C_k(\dots)\in F(\beta_0,\dots,\beta_{k-1})$. Такой корень есть в $\Om$, и он сепарабелен над
$F(\beta_0,\dots,\beta_{k-1})$; по \rf{3.15} и транзитивности сепарабельности (\rf{6.8b}) $L=F(\beta)$ сепарабельно над $F$. Если
$\wp(\beta')=a$, то $\wp(\beta'-\beta)=0$, т.\,е.\ $\beta'=\beta+c$, $c\in W_n(\F_p)$; компоненты $\beta+c$~--- многочлены с
коэффициентами из $\F_p$ от компонент $\beta$ (\rf{5.32l}), так что $F(\beta')\subset F(\beta)$, и по симметрии $F(\beta')=F(\beta)$.
Пусть $\tau\colon L\to\Om$~--- $F$-гомоморфизм. Применяя $\tau$ покомпонентно, получаем $\wp(\tau\beta)=\tau(a)=a$, поэтому
$\tau\beta=\beta+c$ и $\tau(L)=F(\tau\beta)=L$. Каждый корень минимального многочлена над $F$ элемента $x\in L$ имеет вид
$\tau(x)$ для подходящего $\tau$ (гомоморфизм $F[x]\to\Om$ продолжается на $L=F[x][\beta_0,\dots,\beta_{n-1}]$, если
несколько раз применить \rf{2.4}: $\Om$ алгебраически замкнуто), поэтому лежит в $L$. Итак, $L/F$ конечно, сепарабельно и нормально, т.\,е.\ является расширением Галуа.
Для $\sigma\in\Gal(L/F)$ имеем $\sigma\beta-\beta\in\ker\wp=W_n(\F_p)$, и
$\sigma\tau\beta-\beta=\sigma(\tau\beta-\beta)+(\sigma\beta-\beta)=(\tau\beta-\beta)+(\sigma\beta-\beta)$, так как $G$ действует
на $W_n(\F_p)$ тривиально. Если $\sigma\beta=\beta$, то $\sigma$ оставляет на месте образующие $L$, т.\,е.\ $\sigma=1$.

(b) Подгруппа циклической группы $W_n(\F_p)\cong\Z/p^n\Z$ совпадает со всей группой тогда и только тогда, когда она не
содержится в $p\cdot W_n(\F_p)=V\operatorname{Fr}W_n(\F_p)=\{c\mid c_0=0\}$. Нулевая компонента $\sigma\beta-\beta$ равна
$\sigma\beta_0-\beta_0$. Значит, $[L:F]=p^n$ $\iff$ $\sigma\beta_0\ne\beta_0$ для некоторого $\sigma$ $\iff$ $\beta_0\notin F$
$\iff$ $a_0\notin\wp(F)$, так как $\beta_0^p-\beta_0=a_0$ (\rf{5.32c}(b)).

(c) Пусть $G=\Gal(E/F)$ и $\chi\colon G\to W_n(\F_p)$~--- изоморфизм (обе группы циклические порядка $p^n$). Это скрещённый
гомоморфизм в $W_n(E)$, и по \rf{5.32r} $\chi(\sigma)=\sigma\beta-\beta$ для некоторого $\beta\in W_n(E)$. Тогда для
$a=\wp(\beta)$ имеем $\sigma a=\wp(\beta+\chi(\sigma))=\wp(\beta)=a$, т.\,е.\ $a\in W_n(E)^G=W_n(F)$. Элемент $\sigma\in G$,
тождественный на $F(\beta)\subset E$, имеет $\chi(\sigma)=0$, т.\,е.\ $\sigma=1$; значит, $E=F(\beta)$, и $a_0\notin\wp(F)$ по (b).
\end{proof}

\begin{theorem}[5.32t]
Пусть $n\ge1$. Отображение $E\mapsto B(E)=W_n(F)\cap\wp(W_n(E))$ задаёт взаимно однозначное соответствие между конечными
абелевыми расширениями $E\subset\Om$ поля $F$, экспонента группы Галуа которых делит $p^n$, и подгруппами
$\wp(W_n(F))\subset B\subset W_n(F)$ с конечной факторгруппой $B/\wp(W_n(F))$. При этом $[E:F]=(B(E):\wp(W_n(F)))$, и
спаривание $\Gal(E/F)\times B(E)/\wp(W_n(F))\to W_n(\F_p)\cong\Z/p^n\Z$, $(\sigma,b)\mapsto\sigma\beta-\beta$ ($\wp(\beta)=b$),
совершенно. Циклическим расширениям степени $p^n$ отвечают циклические подгруппы порядка $p^n$ группы
$W_n(F)/\wp(W_n(F))$.
\end{theorem}

\begin{proof}
Доказательство дословно повторяет доказательство \rf{5.32d} со следующими заменами: $F$, $\wp(F)$ и $\F_p$ заменяются на
$W_n(F)$, $\wp(W_n(F))$ и $W_n(\F_p)\cong\Z/p^n\Z$; ссылка на \rf{5.32a} заменяется ссылкой на \rf{5.32r}; вместо \rf{5.32c}(a)
используется \rf{5.32s}(a) (все решения $\wp(\beta)=b$ отличаются на элементы $W_n(\F_p)$, и поле $F(\beta)$~--- расширение
Галуа); поле $E_B$ порождается компонентами всех $\beta\in W_n(\Om)$ с $\wp(\beta)\in B$ и является композитом конечного
числа расширений Галуа $F(\beta^{(j)})$ (\rf{3.21}). Вместо линейной алгебры над $\F_p$ используется двойственность конечных
абелевых групп: если экспонента $G$ делит $p^n$, то $|\Hom(G,\Z/p^n\Z)|=|G|$, и для подгруппы $S\subset\Hom(G,\Z/p^n\Z)$
её аннулятор в $G$ имеет порядок $|G|/|S|$. Это \rf{5.30a}(b) и \rf{5.30c}(b) для спаривания
$\Hom(G,\Z/p^n\Z)\times G\to\Z/p^n\Z$, $(\chi,\sigma)\mapsto\chi(\sigma)$ (оно совершенно по \rf{5.30a}(c) и \rf{5.30c}(a)); роль
числа $n$ из \rf{5.30a} и \rf{5.30c} играет здесь $p^n$, а роль $\mu_n$~--- группа $\Z/p^n\Z$: в доказательствах \rf{5.30a} и \rf{5.30c}
используется лишь то, что $\mu_n$~--- циклическая группа порядка $n$. Последнее утверждение: если $E$ циклично степени
$p^n$, то
$B(E)/\wp(W_n(F))\cong\Hom(G,\Z/p^n\Z)\cong\Z/p^n\Z$; обратно, если $B/\wp(W_n(F))$ циклична порядка $p^n$, то
$G\cong\Hom(B/\wp(W_n(F)),\Z/p^n\Z)$ циклична порядка $p^n$.
\end{proof}

\begin{corollary}[5.32u]
В характеристике $p$ каждое циклическое расширение степени $p^n$ ($n\ge1$) содержится в циклическом расширении степени
$p^{n+1}$.
\end{corollary}

\begin{proof}
По \rf{5.32s}(c) $E=F(\beta)$, $\wp(\beta)=a=(a_0,\dots,a_{n-1})$, $a_0\notin\wp(F)$. Положим $a'=(a_0,\dots,a_{n-1},0)\in W_{n+1}(F)$ и
выберем $\beta'\in W_{n+1}(\Om)$ с $\wp(\beta')=a'$. Отбрасывание последней компоненты~--- гомоморфизм колец, перестановочный
с $\operatorname{Fr}$, поэтому для усечения $\beta''=(\beta'_0,\dots,\beta'_{n-1})$ имеем $\wp(\beta'')=a$, и по \rf{5.32s}(a)
$F(\beta'')=F(\beta)=E$. Значит, $E\subset F(\beta')$, а $F(\beta')$ циклично степени $p^{n+1}$ по \rf{5.32s}(b), так как $a'_0=a_0\notin\wp(F)$.
\end{proof}

\begin{remark}[5.32v]
В характеристике $0$ это неверно: $\Q(i)/\Q$ не содержится ни в каком циклическом расширении степени $4$. Действительно,
пусть $K\supset\Q(i)$, $K/\Q$ циклично степени $4$ с образующей $\sigma$. Тогда $K=\Q(i)(\sqrt w)$, $w\in\Q(i)$;
$\sigma^2$~--- нетривиальный элемент $\Gal(K/\Q(i))$, так что $\sigma^2\sqrt w=-\sqrt w$, а $\sigma|_{\Q(i)}$~--- комплексное
сопряжение, так что $(\sigma\sqrt w)^2=\bar w$. Элемент $\theta=\sqrt w\cdot\sigma\sqrt w$ удовлетворяет
$\sigma\theta=\sigma\sqrt w\cdot\sigma^2\sqrt w=-\theta$, поэтому $\theta\notin\Q$, $\sigma^2\theta=\theta$, т.\,е.\ $\theta\in\Q(i)$, и
$\theta^2=w\bar w=|w|^2>0$. Но если $\theta=x+yi\in\Q(i)\smallsetminus\Q$ и $\theta^2=x^2-y^2+2xyi\in\Q$, то $x=0$ и
$\theta^2=-y^2<0$~--- противоречие.
\end{remark}

\begin{example}[5.32w: циклическое расширение степени $4$ над $\F_2$]
Пусть $p=2$, $n=2$, $a=(1,0)\in W_2(\F_2)$; $a_0=1\notin\wp(\F_2)=\{0\}$, и по \rf{5.32s} поле $\F_2(\beta_0,\beta_1)$, где
$\wp(\beta)=a$, циклично степени $4$, т.\,е.\ это $\F_{16}$. Запишем уравнение $\operatorname{Fr}(\beta)=\beta+a$ по
компонентам с помощью $S_1=X_1+Y_1-X_0Y_0$ (в характеристике $2$ знак не важен):
\[
\beta_0^2=\beta_0+1,\qquad\beta_1^2=\beta_1+0+\beta_0\cdot1,\qquad\text{т.\,е.}\qquad\beta_0^2+\beta_0+1=0,\quad\beta_1^2+\beta_1=\beta_0.
\]
Первое уравнение даёт $\F_4=\F_2(\beta_0)$. Второе~--- уравнение Артина--Шрайера над $\F_4$; оно неприводимо, так как
$\operatorname{Tr}_{\F_4/\F_2}(\beta_0)=\beta_0+\beta_0^2=1\ne0$ (\rf{5.32f}). Исключая $\beta_0$, получаем
$(\beta_1^2+\beta_1)^2+(\beta_1^2+\beta_1)+1=\beta_1^4+\beta_1+1=0$: $\F_{16}=\F_2[X]/(X^4+X+1)$, $\beta_1=X$,
$\beta_0=X^2+X$. Образующая $\phi(x)=x^2$ группы $\Gal(\F_{16}/\F_2)$ действует по правилу $\phi(\beta)=\operatorname{Fr}(\beta)=\beta+a$,
т.\,е.\ соответствует $1\in\Z/4\Z=W_2(\F_2)$; $\phi^2(\beta)=\beta+2\cdot a=\beta+(0,1)=(\beta_0,\beta_1+1)$, и действительно
$\beta_1^4=\beta_1+1$. Подполе $\F_4$ задаётся усечением $a\mapsto a_0=1$.

\begin{center}
\begin{tikzpicture}[every node/.style={font=\small}]
\node (A) at (0,0) {$\F_2$};
\node at (1.3,0) {$\subset$};
\node (B) at (3.0,0) {$\F_4=\F_2(\beta_0)$};
\node at (4.7,0) {$\subset$};
\node (C) at (6.6,0) {$\F_{16}=\F_4(\beta_1)$};
\draw[decorate,decoration={brace,mirror,amplitude=5pt},cold] (-0.3,-0.35) -- node[below=6pt,align=center]{$\beta_0^2+\beta_0=1$\\ $\Z/2$} (3.0,-0.35);
\draw[decorate,decoration={brace,mirror,amplitude=5pt},cold] (3.1,-0.35) -- node[below=6pt,align=center]{$\beta_1^2+\beta_1=\beta_0$\\ $\Z/2$} (7.8,-0.35);
\draw[decorate,decoration={brace,amplitude=5pt},cnew] (-0.3,0.35) -- node[above=6pt]{$\Z/4$: \ $\phi(\beta)=\beta+(1,0)$ в $W_2(\F_{16})$} (7.8,0.35);
\end{tikzpicture}
\end{center}
\figcap{Каждый из двух шагов башни~--- расширение Артина--Шрайера с группой $\Z/2$; то, что всё расширение
циклично (группа $\Z/4$, а не $(\Z/2)^2$), видно из одного уравнения в векторах Витта длины $2$.}

Так же при $p=3$, $a=(1,0)\in W_2(\F_3)$: $S_1=X_1+Y_1-X_0^2Y_0-X_0Y_0^2$ даёт $\beta_0^3-\beta_0=1$,
$\beta_1^3-\beta_1=-\beta_0^2-\beta_0$, и $\beta_1$~--- корень неприводимого над $\F_3$ многочлена
$X^9-X^6-X^4-X^3-X^2+1$ (результант по $\beta_0$), т.\,е.\ $\F_3(\beta_0,\beta_1)=\F_{3^9}$, циклическое расширение степени $9$.
\end{example}

\begin{remark}[5.32x: абелевы расширения экспоненты $n$ при $p\mid n$]
Пусть $\operatorname{char}F=p$, $n=p^km$, $p\nmid m$, и $E/F$~--- конечное абелево расширение, экспонента группы $G$
которого делит $n$ (ср.\ \rf{5.31j}). Тогда $G=G_p\times G'$, где $G_p$~--- $p$-примарная компонента (экспонента делит $p^k$), а $|G'|$ не
делится на $p$ (экспонента делит $m$). Положим $E_1=E^{G'}$, $E_2=E^{G_p}$. По основной теореме (\rf{3.17}) $E_1/F$ и
$E_2/F$~--- абелевы расширения с группами $G/G'\cong G_p$ и $G/G_p\cong G'$, $E_1\cap E_2=E^{G}=F$ и $E_1E_2=E^{G_p\cap G'}=E$.
Обратно, композит абелевых расширений экспонент, делящих $p^k$ и $m$, абелев экспоненты, делящей $n$ (\rf{3.21}). Итак,
\begin{itemize}
\item \emph{$p$-часть} $E_1$ описывается теорией Артина--Шрайера--Витта (\rf{5.32t}) с $W_k(F)$;
\item \emph{часть, взаимно простая с $p$,} описывается теорией Куммера: если $\mu_m\subset F$, то по \rf{5.30} $E_2=F[B^{1/m}]$,
$B\subset F^\times$; в общем случае $E_2F'/F'$, где $F'=F[\zeta_m]$, абелево экспоненты, делящей $m$ (\rf{3.19}), и к нему
применима \rf{5.30} над $F'$.
\end{itemize}
\end{remark}

\begin{center}
\small
\begin{tabular}{@{}p{0.17\textwidth}p{0.23\textwidth}p{0.23\textwidth}p{0.27\textwidth}@{}}
\toprule
 & Куммер & Артин--Шрайер & Артин--Шрайер--Витт\\
\midrule
условия & $\mu_n\subset F$, $\operatorname{char}F\nmid n$;\newline экспонента делит $n$ & $\operatorname{char}F=p$;\newline экспонента $p$ & $\operatorname{char}F=p$;\newline экспонента делит $p^n$\\
уравнение & $\alpha^n=a$ & $\alpha^p-\alpha=a$ & $\operatorname{Fr}(\beta)-\beta=a$ в $W_n$\\
образующая & $\sigma\alpha=\zeta\alpha$ & $\sigma\alpha=\alpha+1$ & $\sigma\beta=\beta+1$, $1=(1,0,\dots,0)$\\
ядро & $\mu_n\cong\Z/n$ & $\F_p\cong\Z/p$ & $W_n(\F_p)\cong\Z/p^n$\\
параметры & $F^\times/F^{\times n}$ & $F/\wp(F)$ & $W_n(F)/\wp(W_n(F))$\\
Гильберт~90 & $H^1(G,E^\times)=0$ (\rf{5.23}) & $H^1(G,E)=0$ (\rf{5.32a}) & $H^1(G,W_n(E))=0$ (\rf{5.32r})\\
классификация & \rf{5.27}, \rf{5.28}, \rf{5.30} & \rf{5.32c}, \rf{5.32d} & \rf{5.32s}, \rf{5.32t}\\
\bottomrule
\end{tabular}
\end{center}
\figcap{Три теории абелевых расширений. Во всех трёх случаях классифицирующая группа~--- коядро гомоморфизма
($x\mapsto x^n$, $\wp$, $\operatorname{Fr}-1$), ядро которого~--- циклическая группа значений спаривания, а сюръективность
на когомологиях обеспечивает соответствующая теорема Гильберта~90.}
